\begin{discussion}

	研究流行病学仓室模型时，流行病学调查、微分方程方法和神经网络方法是三种相互关联且互补的工具。首先，流行病学调查在这个过程中扮演着基础性的角色。
	它提供了必要的数据基础，这些数据可以用来校准模型参数，验证模型的预测，以及作为输入数据供模型使用。流行病学调查的数据不仅丰富，而且具有真实世界的背景，使得模型更有实际意义。

	与此同时，微分方程方法和神经网络方法相互互补，对流行病学仓室模型的研究提供了不同的视角。微分方程方法建立在流行病学原理之上，提供了基于传统流行病学知识的模型。这些模型能够捕捉疾病传播的动态变化，为疾病传播的数学建模提供了强大的工具。
	另一方面，神经网络方法处理复杂的非线性关系，特别适用于疾病传播过程中的随机性和复杂性。神经网络可以通过学习数据中的模式和趋势来改进模型的预测性能，尤其在考虑到多种因素时，其表现出色。

	将这些方法综合使用，有助于提高对流行病学过程的理解和预测能力。这不仅有助于更好地理解疾病的传播方式和趋势，还能为公共卫生政策制定和干预措施的制定提供更有根据的指导。
	通过综合运用这些工具，可以更有效地应对流行病威胁，保护公众健康。
\end{discussion}